\subsection{Tree responses survive thick equitable cell quotients}
\label{subsec:equitable-tree-quotients}

The tree-response proof does not require the original graph itself to be a
tree.  It requires only a tree after exact symmetry compression.  Let
\begin{equation}
 \mathcal P=\{C_a:a\in\mathcal A\}
 \label{eq:equitable-cell-partition}
\end{equation}
be a partition of $V$ with a root cell $C_r=\{o\}$.  Call it an
\emph{equitable tree partition} if every $v\in C_a$ has a common degree
$d_a$ and, for every cell $C_b$, a common count
\begin{equation}
 \nu_{ab}:=|N(v)\cap C_b|,
 \label{eq:equitable-cell-count}
\end{equation}
and if the graph on distinct cell indices with
$a\sim b$ when $\nu_{ab}>0$ is a tree rooted at $r$.  Edges within one cell
are allowed.  Put $m_a:=|C_a|$.  As before,
$m_a\nu_{ab}=m_b\nu_{ba}$.

The partition is \emph{locally enumerable} if an exposed vertex label carries
its cell identifier at $O(1)$ cost.  Scanning an active cell then groups its
incident labels by neighboring cell.  This representation supplies neither
the active cells nor their depth.  Moreover, it is locally auditable: all
degree and incidence counts of an active cell are read during its one scan.

\begin{theorem}[Exact RPPR at the product scale on equitable tree quotients]
\label{thm:equitable-tree-quotient}
Suppose a single-seed instance has a locally enumerable equitable tree
partition.  With
\begin{equation}
 e_a:=\frac{\mathbf 1_{C_a}}{\sqrt{m_a}},
 \label{eq:equitable-cell-basis}
\end{equation}
the restriction $\widehat Q$ of $Q$ to the cell-constant subspace is a
symmetric positive-definite Stieltjes matrix supported on the quotient tree,
with
\begin{align}
 \widehat Q_{aa}
 &=\frac{1+\alpha}{2}
   -\frac{1-\alpha}{2}\frac{\nu_{aa}}{d_a},
 \label{eq:equitable-tree-quotient-diagonal}\\
 \widehat Q_{ab}
 &=-\frac{1-\alpha}{2}
   \sqrt{\frac{\nu_{ab}\nu_{ba}}{d_ad_b}}
 \quad(a\sim b).
 \label{eq:equitable-tree-quotient-edge}
\end{align}
Its demand is
\begin{equation}
 \widehat c_r
 =\frac{\alpha}{\sqrt{d_r}}-\alpha\rho\sqrt{d_r},
 \qquad
 \widehat c_a=-\alpha\rho\sqrt{m_ad_a}
 \quad(a\neq r).
 \label{eq:equitable-tree-quotient-demand}
\end{equation}
The RPPR optimum is cell-constant.  Writing
\begin{equation*}
 \mathcal A^\star:=\{a:C_a\subseteq S^\star(\rho)\},
\end{equation*}
there is an exact adjacency-local algorithm that discovers these cells and
returns the full optimum using
\begin{align}
 \Work_{\rm cell\text{-}tree}
 &\leq
 O\!\left(
  \vol(S^\star)
  +\log(1+\Delta_{\mathcal P})
   \sum_{a\in\mathcal A^\star}
   \operatorname{depth}_{\mathcal P}(a)
 \right)
 \notag\\
 &\leq
 O\!\left(
  \vol(S^\star)
  [1+R_\star\log(1+\Delta_{\mathcal P})]
 \right)
 =\widetilde O\!\left(\frac1{\rho\sqrt\alpha}\right),
 \label{eq:equitable-tree-quotient-work}
\end{align}
where $\Delta_{\mathcal P}$ is the maximum quotient-tree degree.  It uses
$O(\vol(S^\star))$ storage and no support, radius, or confinement oracle.
Every degree and incidence identity encountered on the active exploration is
audited.  An observed failure hands the current safe region and its exact
coordinatewise violating batch to the general boundary gate, without
compromising correctness.
\end{theorem}

\begin{proof}
Let $P_{\mathcal P}$ be orthogonal averaging onto the cell-constant subspace.
The degree and incidence identities make this subspace invariant under
$D^{-1/2}AD^{-1/2}$.  Symmetry makes its orthogonal complement invariant as
well, so $P_{\mathcal P}Q=QP_{\mathcal P}$.  The shifted demand is constant
on each cell.  Therefore, for $z=(I-P_{\mathcal P})x$,
\begin{equation*}
 \Phi_\rho(x)
 =\Phi_\rho(P_{\mathcal P}x)+\frac12z^\top Qz.
\end{equation*}
Orthogonal averaging preserves nonnegativity, and positive definiteness puts
the unique optimum in the cell-constant subspace.  Direct projection, using
edge balance, gives
\cref{eq:equitable-tree-quotient-diagonal,eq:equitable-tree-quotient-edge,eq:equitable-tree-quotient-demand}.
As a restriction of $Q$, the quotient has spectrum in $[\alpha,1]$.

Root the quotient tree at $r$.  The response construction in
\cref{thm:tree-response,cor:inverse-tree-response} uses only tree sparsity,
the Stieltjes signs, strictly negative nonroot demands, and the spectral lower
bound $\alpha I$.  It therefore applies verbatim to $(\widehat Q,\widehat c)$.
For a child cell $a$ with parent $p(a)$, its first event is
\begin{equation}
 \widehat\tau_a
 =\frac{\alpha\rho\sqrt{m_ad_a}}
        {-\widehat Q_{p(a),a}}.
 \label{eq:equitable-cell-first-event}
\end{equation}
Scanning all adjacency lists in the active parent cell enumerates every label
of $C_a$: equitability gives each such label the positive common number
$\nu_{a,p(a)}$ of parent neighbors.  The same scan gives $m_a$,
$\nu_{p(a),a}$, and $\nu_{a,p(a)}$, while degree queries give $d_a$.
Thus \cref{eq:equitable-cell-first-event} is known before $C_a$ is scanned.
Only if that event is consumed does the algorithm scan $C_a$, learn its
diagonal and child couplings, and initialize its child-event heap.

Consequently every active original adjacency list is scanned once, for
$O(\vol(S^\star))$ work.  There is one activation label per active cell, and
the lazy merge advances it once through each quotient ancestor, giving the
first line of \cref{eq:equitable-tree-quotient-work}.  Because the quotient
graph is a rooted tree and every child-cell vertex has a parent neighbor,
$\operatorname{depth}_{\mathcal P}(a)$ equals the original root distance of
the vertices in $C_a$.  Hence it is at most $R_\star$ for an active cell.
Also $|\mathcal A^\star|\leq|S^\star|\leq\vol(S^\star)$.  This proves the
second line; the support-volume bound and
\cref{lem:graph-support-radius} give the final order.

Top-down evaluation of the quotient responses returns $y_a$ and then
$x_v=y_a/\sqrt{m_a}$ for $v\in C_a$.  This is the full optimum by the
orthogonal splitting.  Finally, before a child event is consumed, the parent
scan checks the degrees and backward incidences of all its exposed labels; if
they differ, their individual violations define the safe gate batch.  After a
whole child cell is admitted, its own scan checks the internal and forward
incidences.  A failure at that point follows a boundary-gate violation that
has already admitted the cell safely.  The same Schur-complement argument as
in \cref{thm:certifying-shell-fast-path} permits an exact fallback in either
case.
\end{proof}

This theorem strictly extends the one-cell-per-shell result.  Unequal
clique, independent-set, or biregular blow-ups of a rooted tree may have
several inequivalent thick cells in one distance shell, unbounded treewidth,
and arbitrarily large biconnected cores.  Their original graph need not be
root-distance equitable, but the exact quotient still has tree responses and
meets the intended product scale.

The cell labels in the preceding statement are convenient but not necessary.
The reason is causal: before an inactive cell activates, only its degree and
its incidences with already scanned cells affect its obstacle slack.  Those
quantities are visible from the scanned side of the boundary.  Cells that
cannot yet be distinguished therefore have the same next event and may be
opened together.  Their different futures can be separated after, but not
before, that event.

For a scanned set $U$ containing $o$, form a emph{visible refinement}
$\mathcal R_U$ as follows.  Initially color vertices of $U$ by the root flag
and their full graph degree.  Repeatedly split a color class whenever two of
its vertices have different neighbor counts into some current class.  In an
online run, colors present before an expansion are retained as base colors,
so later refinement may split but never merge old classes.  All incidences
used here are known because every adjacency list in $U$ has been scanned.
For an exposed label $v\in N(U)\setminus U$, define its visible signature
by
\begin{equation}
 \sigma_U(v)
 :=\left(
   d_v,
   \bigl(|N(v)\cap B|\bigr)_{B\in\mathcal R_U}
 \right).
 \label{eq:visible-boundary-signature}
\end{equation}
The counts are accumulated from the scanned endpoints in $U$; the adjacency
list of $v$ is not read.

\begin{lemma}[Hidden tree cells refine causally]
\label{lem:hidden-cell-refinement}
Suppose $G$ admits an equitable tree partition $\mathcal P$ with root cell
$\{o\}$, but $\mathcal P$ is not supplied to the algorithm.  Let $U$ be a
union of cells whose quotient indices form a rooted connected subtree.  Then:
\begin{enumerate}[label=(\roman*),leftmargin=*]
 \item every class of $\mathcal R_U$ is a union of cells of $\mathcal P$,
 and the distinct-class interaction graph of $\mathcal R_U$ is a rooted
 tree;
 \item the restricted obstacle solution on $U$, and every affine piece of
 the root-conditioned homotopy on $U$, is constant on the classes of
 $\mathcal R_U$;
 \item every signature class in $N(U)\setminus U$ is a union of whole
 boundary cells of $\mathcal P$, and all its labels have the same affine
 obstacle slack;
 \item if such a signature class reaches the next event, admitting the whole
 class and scanning its lists is exact.  Refining afterward preserves the
 common solution value at the event and forks only the future affine pieces.
\end{enumerate}
\end{lemma}

\begin{proof}
The initial colors do not split a cell of $\mathcal P$: the root cell is a
singleton and degree is constant inside every other cell.  If each current
color is a union of $\mathcal P$-cells, equitability makes the neighbor count
into that color constant on every $\mathcal P$-cell.  Induction over the
refinement steps therefore proves the first assertion of (i).

The unique root color propagates through one quotient edge per refinement
round, so two visible classes at different quotient depths cannot merge.
Moreover, if two nonroot $\mathcal P$-cells remain in one visible class, their
unique parent cells must lie in one visible class: otherwise their count
vectors into the two parent classes differ.  Consequently each nonroot
visible class has one visible parent, proving that its interaction graph is a
rooted tree.

Stability of the refinement says that degree and all within-$U$ incidence
counts are constant on a visible class.  Hence its cell-constant subspace is
invariant under $Q_{UU}$, while the restricted shifted demand is constant on
the same classes.  Orthogonal averaging, exactly as in the proof of
\cref{thm:equitable-tree-quotient}, puts the restricted obstacle minimizer and
each fixed-support homotopy solution in that subspace.  This proves (ii).

Let $C_a$ be a quotient cell adjacent to $U$.  Its unique parent cell belongs
to $U$, and $\nu_{a,p(a)}>0$.  Scanning that parent therefore exposes every
label of $C_a$.  Equitability makes
\cref{eq:visible-boundary-signature} constant throughout $C_a$, so a
signature class is a union of whole boundary cells.  If $x_U(t)$ is one
affine homotopy piece, the inactive slack at $v$ is
\begin{equation*}
 g_v(t)=(Q_{vU}x_U(t))-c_{\rho,v}.
\end{equation*}
By (ii), the visible signature determines every term on the right.  Thus the
slack is common and affine on a signature class, proving (iii).

At its zero, the entire tied class satisfies the activation equality.  Before
the zero all its coordinates are zero, so unrevealed internal or forward
edges have made no contribution.  The coordinates are continuous at the
event.  Carrying the old colors as base colors and scanning the new bundle
can only split classes of the old exact quotient; every new class inherits
the common event value, after which its own affine continuation is computed
from the refined quotient.  This proves (iv).
\end{proof}

\begin{lemma}[Amortized visible refinement and response forks]
\label{lem:visible-refinement-work}
Let $k_\star$ be the number of active cells in a witnessing equitable tree
partition.  Along the exact homotopy up to its root zero, visible refinement
can be maintained in
\begin{equation}
 O\!\left(
  \vol(S^\star)\log(1+\vol(S^\star))
 \right)
 \label{eq:visible-refinement-work}
\end{equation}
total work and $O(\vol(S^\star))$ storage.  The additional lazy-response work
created by class splits is at most a constant multiple of
\begin{equation}
 \log(1+\Delta_{\mathcal P})
 \sum_{a\in\mathcal A^\star}
 \operatorname{depth}_{\mathcal P}(a).
 \label{eq:visible-response-fork-work}
\end{equation}
\end{lemma}

\begin{proof}
Store every incidence when its active endpoint is first scanned.  Maintain
the stable colors with a splitter queue, always retaining one largest part of
a split under its old record and placing the other parts on the queue.  An
incidence is revisited only when one of its endpoint classes is a newly
inserted class or is charged to a part of at most half the former size.  The
first case occurs once and the second at most
$\log_2(1+\vol(S^\star))$ times.  Boundary-label counts are updated by the
same stored incidences.  This is the standard smaller-half refinement charge
and proves \cref{eq:visible-refinement-work}.

Past affine pieces are immutable.  When a visible class splits, retain their
shared prefix and create records only for the new future pieces.  Each split
strictly increases the active visible-class count, and no true cell is split
by \cref{lem:hidden-cell-refinement}; hence the total number of created active
records is $O(k_\star)$.  More precisely, the class histories form refinement
forests with no unary internal nodes, hence fewer than twice as many records
as final visible classes.  Every record and all of its descendant classes
have the same quotient depth.  Their total propagation depth is therefore at
most twice the sum of the witnessing active-cell depths.  A new record changes
only the response of its unique visible parent, so the lazy merge propagates
it once through each ancestor.  Heap maintenance supplies the logarithm in
\cref{eq:visible-response-fork-work}, exactly as in
\cref{thm:lazy-tree-response}.
\end{proof}

\begin{theorem}[Representation-free product bound for hidden equitable trees]
\label{thm:hidden-equitable-tree-quotient}
Suppose a single-seed graph admits an equitable tree partition with root cell
$\{o\}$.  The partition and its cell identifiers need not be supplied.
Ordinary degree queries and adjacency-list scans suffice to discover and
solve exact RPPR using
\begin{align}
 \Work_{\rm hidden\text{-}cell\text{-}tree}
 &=O\!\left(
   \vol(S^\star)\log(1+\vol(S^\star))
   +\log(1+\Delta_{\mathcal P})
    \sum_{a\in\mathcal A^\star}
    \operatorname{depth}_{\mathcal P}(a)
  \right)
 \notag\\
 &=\widetilde O\!\left(
   \vol(S^\star)(1+R_\star)
  \right)
 =\widetilde O\!\left(\frac1{\rho\sqrt\alpha}\right).
 \label{eq:hidden-equitable-tree-work}
\end{align}
The algorithm uses $O(\vol(S^\star))$ storage and no support, radius,
confinement, partition, or cell-identifier oracle.

On an arbitrary graph the same procedure is certifying.  If a revealed
refined quotient is not a tree, it hands the exact active region and the
individual tied violation batch to the graph-universal homotopy or boundary
gate.  Thus the hidden-tree property is used only for the work bound, not for
correctness.
\end{theorem}

\begin{proof}
Run the lazy quotient-response algorithm of
\cref{thm:equitable-tree-quotient}, replacing supplied cells by the visible
classes and boundary signatures above.  By
\cref{lem:hidden-cell-refinement}, its current quotient and every event are
exact; a bundle is scanned only after its common event is consumed.  The
root iterator therefore stops at the unique root zero and exposes exactly
$S^\star$.  The original adjacency lists of active vertices are scanned once.
Add the refinement and fork charges of
\cref{lem:visible-refinement-work} to the ordinary lazy-response work.  Since
$|\mathcal A^\star|\leq\vol(S^\star)$ and every active quotient depth is at
most $R_\star$, the second line of
\cref{eq:hidden-equitable-tree-work} follows.  The support-volume and radius
bounds give the last equality.

For a graph outside the promised class, visible refinement still gives an
exact quotient of the scanned principal problem, and every consumed boundary
label has an exact zero slack.  The first non-tree interaction is therefore
an observed structural witness, not a speculative failure.  Monotonicity of
the graph-universal homotopy makes the active set a subset of $S^\star$.
Re-solve that scanned set at the fixed target $\rho$ and
\cref{thm:condition-free-graph} completes an exact fallback.
\end{proof}

Equivalently, the useful structural condition is intrinsic: the stable
root-and-degree equitable refinement may have a tree interaction graph.  The
algorithm discovers only the portion that becomes active.  A precomputed or
trusted symmetry partition is not part of the access model.
